Let $A$ be an abelian group written additively. By an endomorphism of $A$ we mean a group homomorphism $f: A \to A$; that is,

$$f(a + b) = f(a) + f(b) \quad (a, b \in A).$$

The set $E$ of all such endomorphisms of $A$ forms an abelian group with respect to the addition $(f, g) \mapsto f + g$ defined by

$$(f + g)(a) = f(a) + g(a) \quad (a \in A).$$

The identity and the inverse (negative) are given by

$$0(a) = 0 \quad \text{and} \quad (-f)(a) = -f(a) \quad (a \in A).$$

With composition as multiplication, $E$ becomes a ring, called the endomorphism ring of $A$.
